% =====================================================================
%  Literature audit: results in conflict with Douna et al. (2018)
%  Generated 2026-09-08.  Compile with: pdflatex contradictions_douna2018.tex
% =====================================================================
\documentclass[11pt,a4paper]{article}

\usepackage[T1]{fontenc}
\usepackage[utf8]{inputenc}
\usepackage[margin=2.4cm]{geometry}
\usepackage{amsmath,amssymb}
\usepackage{booktabs}
\usepackage{longtable}
\usepackage{enumitem}
\usepackage{xcolor}
\usepackage[colorlinks=true,linkcolor=blue!50!black,urlcolor=blue!50!black,
            citecolor=blue!50!black]{hyperref}

\newcommand{\HYP}{\textcolor{blue!60!black}{\textbf{HYPOTHESIS}}}
\newcommand{\RES}{\textcolor{red!60!black}{\textbf{RESULTS}}}
\newcommand{\BOTH}{\textcolor{purple!70!black}{\textbf{HYPOTHESIS + RESULTS}}}
\newcommand{\flag}[1]{\textcolor{orange!75!black}{\textsf{\small [#1]}}}

\title{\bfseries Papers whose results conflict with\\
Douna et al. (2018), \emph{Microquasars as heating sources of the\\
intergalactic medium during reionization of the Universe}\\[6pt]
\large A provenance-tracked literature audit}
\author{Prepared for L.\ Carvalho}
\date{8 September 2026}

\begin{document}
\maketitle

% =====================================================================
\section{Scope, method and audit trail}
% =====================================================================

\subsection{Target paper}

\begin{quote}
\textbf{Douna, V.~M., Pellizza, L.~J., Laurent, P.\ \& Mirabel, I.~F.\ (2018)},
\emph{Microquasars as heating sources of the intergalactic medium during
reionization of the Universe}, MNRAS \textbf{474}, 3488--3499.\\
DOI: \href{https://doi.org/10.1093/mnras/stx2983}{10.1093/mnras/stx2983} ---
arXiv:\href{https://arxiv.org/abs/1711.07374}{1711.07374} ---
bibcode \texttt{2018MNRAS.474.3488D}.\\
Local file: \texttt{Douna (2018) - Microquasars as heating sources\dots}\ (parent folder).
\end{quote}

\subsection{How the conflicting set was assembled}

Four steps, all executed on 8 September 2026:

\begin{enumerate}[leftmargin=*,itemsep=2pt]
\item \textbf{Full read of the target paper.} Sections~1--6 of the published
      MNRAS version were read in full and reduced to the list of testable
      claims in Section~\ref{sec:claims}.
\item \textbf{In-folder search.} All papers in the parent folder were
      text-searched for statements in conflict with those claims.
      Findings in Section~\ref{sec:infolder}.
\item \textbf{Web search --- arXiv.} The arXiv API was queried across
      \texttt{astro-ph.HE}, \texttt{astro-ph.CO} and \texttt{astro-ph.GA} on
      the axes: cosmic-ray heating/ionization of the IGM; cosmic-ray transport
      and magnetic confinement at high $z$; X-ray-binary contribution to
      reionization; the metallicity--$L_{\rm X}$ scaling; 21-cm observational
      constraints on IGM temperature. Missing bibliographic data were resolved
      through Crossref.
\item \textbf{Web search --- NASA ADS.} The ADS API was queried with a
      user-supplied token for (a) the authoritative citation list
      \texttt{citations(bibcode:2018MNRAS.474.3488D)}, (b) full-text searches
      (\texttt{full:"Douna"}, \texttt{full:"microquasar" full:"reionization"}),
      and (c) refereed topical searches on the same axes as step~3.
\item \textbf{Screening.} 19 candidates were downloaded; 15 were retained in
      \texttt{against/} and 4 were moved to \texttt{screened\_no\_conflict/}
      after full-text inspection failed to yield a conflicting statement
      (Section~\ref{sec:screened}).
\end{enumerate}

\subsection{What the ADS query changed (recall audit)}
\label{sec:adsdelta}

The ADS cross-check materially revised two things.

\begin{enumerate}[leftmargin=*,itemsep=3pt]
\item \textbf{The citation count was wrong. ADS reports 18 citing papers, not
      14.} Semantic Scholar, OpenAlex, OpenCitations and INSPIRE-HEP jointly
      missed four items, all \emph{Bolet\'in de la Asociaci\'on Argentina de
      Astronom\'ia} proceedings with no DOI and no arXiv posting, which is why
      no Crossref-derived index carries them:
      \begin{itemize}[leftmargin=1.2cm,itemsep=1pt]
      \item Escobar, Pellizza \& Romero 2021, BAAA 62, 265,
            \emph{Microcu\'asares como fuentes de rayos c\'osmicos}
            (\texttt{2021BAAA...62..265E});
      \item Garate N\'u\~nez, Escobar, Pellizza \& Bosch-Ramon 2021, BAAA 62, 234,
            \emph{Reionizaci\'on por jets de n\'ucleos gal\'acticos activos}
            (\texttt{2021BAAA...62..234G});
      \item Carvalho, Escobar \& Pellizza 2024, BAAA 65, 240,
            \emph{Cosmic rays at the epoch of reionization}
            (\texttt{2024BAAA...65..240C});
      \item Pellizza, Badaracco, Carvalho, Escobar \& Bignone 2025, BAAA 66, 377,
            \emph{X-ray binary populations in local galaxies}
            (\texttt{2025BAAA...66..377P}).
      \end{itemize}
      All four are now downloaded into the parent folder. \flag{note} The last
      two are your own group's work with L.~J.\ Pellizza (Douna's co-author);
      they continue the line rather than conflict with it, and none of the four
      contradicts any claim in Section~\ref{sec:claims}.
\item \textbf{Two new conflicting papers were found}, one of which is now the
      strongest entry in the catalogue (Dhandha et al.\ 2025, \S3.1) and one
      moderate (Kaur et al.\ 2022, \S3.11). ADS full-text search also surfaced
      three papers that turned out to \emph{support} Douna et al.'s premises
      and were reclassified (Section~\ref{sec:screened}).
\end{enumerate}

\noindent\flag{recall, after the cross-check} The four open indices had
\emph{complete} recall on refereed, DOI-bearing literature and \emph{zero}
recall on Spanish-language national proceedings. Anything comparable in, e.g.,
RevMexAA or IAU regional proceedings would have been missed before this step.

\subsection{Reading conventions used in this document}

\begin{itemize}[leftmargin=*,itemsep=2pt]
\item Every entry is tagged \HYP{} when the conflict is with material in
      Douna et al.'s Introduction (\S1) or Model/Simulations (\S2), and
      \RES{} when it is with material in their Results (\S3--5) or
      Discussion and Conclusions (\S6).
\item \textbf{Excerpts are verbatim} from the PDF stored in \texttt{against/}.
      Text was extracted with \texttt{pdftotext}; mathematical symbols that the
      extractor mangled have been restored to their printed form and every such
      restoration is flagged in the entry. No wording has been altered.
\item \textbf{Conflict strength} is graded
      \textsc{direct} (the paper's own stated result is logically incompatible
      with a Douna claim),
      \textsc{moderate} (incompatible under stated assumptions, or contradicts
      an input rather than an output), and
      \textsc{indirect} (undercuts a premise or the motivating context without
      addressing the claim itself).
\end{itemize}

\flag{Confidence, stated up front} \textbf{None of the 15 papers in
\texttt{against/} cites Douna et al.\ (2018), and none is written as a rebuttal
of it.} Every conflict below is a conflict I have identified between
independently obtained results; it is not a disagreement the authors themselves
declare. This is the single most important caveat in this document and it
applies to every entry. See Section~\ref{sec:global} for the corollary.

% =====================================================================
\section{The claims of Douna et al.\ (2018) under test}
\label{sec:claims}
% =====================================================================

\subsection{Hypotheses, premises and modelling choices (\S1--\S2)}

\begin{description}[leftmargin=1.6cm,style=nextline,itemsep=3pt]
\item[\textbf{H1}] X-ray binaries, and hence microquasars (MQs), were more
  numerous and more luminous at high $z$ because of the low metallicity of
  their parent populations (\S1, citing Mirabel et al.\ 2011; Fragos et al.\
  2013; Douna et al.\ 2015; Brorby et al.\ 2016).
\item[\textbf{H2}] The X-ray channel is insufficient: ``recent works suggest
  that the contribution of the X-rays emitted by these systems to heating and
  ionization of the IGM during the EoR is marginal at most (Madau \& Fragos
  2017)'' (\S1). This is the stated motivation for turning to CR electrons.
\item[\textbf{H3}] CR \emph{electrons} are the carrier worth modelling first,
  ``as these particles are expected to make the largest contribution (Tueros
  et al.\ 2014)'' (\S1, final paragraph).
\item[\textbf{H4}] The reference (``warm'') scenario assumes a pre-heated IGM at
  $T_{\rm IGM}\sim10^{4}$~K at $z=10$, which fixes the case-B recombination
  coefficient $\alpha_{\rm B}\approx2.6\times10^{-13}\,{\rm cm^{3}\,s^{-1}}$
  used throughout \S5.1.
\item[\textbf{H5}] The source is a single MQ (or population) with a power-law
  electron spectrum, index $\alpha\in[2,3]$ (fiducial 2.5), total jet kinetic
  luminosity $L_{\rm k}=10^{40}\,{\rm erg\,s^{-1}}$, cut-offs free in
  10~keV--1~PeV (\S5).
\item[\textbf{H6}] \textbf{Magnetic fields are omitted.} ``there are still many
  uncertainties in the knowledge of these fields in the early Universe\dots
  Therefore, we will treat magnetic fields separately in a companion paper''
  (\S2.1). Consequently there is no synchrotron cooling and no diffusive
  confinement; transport is rectilinear.
\item[\textbf{H7}] EBL, photoionization by secondary photons and galactic
  photon fields are omitted; the authors argue each omission is
  \emph{conservative} (i.e.\ their rates are lower bounds) (\S2.1, \S6).
\item[\textbf{H8}] ISM and IGM are cold, homogeneous, partially ionized
  primordial plasmas; $n_{\rm ISM}=1\,{\rm cm^{-3}}$,
  $n_{\rm IGM}=2.4\times10^{-4}\,{\rm cm^{-3}}$; cosmological redshifting of
  the CMB during propagation is neglected (\S2.1, \S2.3).
\end{description}

\subsection{Results and conclusions (\S3--\S6)}

\begin{description}[leftmargin=1.6cm,style=nextline,itemsep=3pt]
\item[\textbf{R1}] Low-energy electrons (up to several tens of keV), which carry
  most of the ionizing power, cannot escape primeval galaxies; any effective
  mechanism must be two-stage (\S6, point i).
\item[\textbf{R2}] The escape fraction depends on energy \emph{and} on source
  spectrum and may exceed unity, because cooling creates new electrons
  (\S6, point ii).
\item[\textbf{R3}] Under the warm scenario, MQs alone can maintain
  $f_{\rm ion}^{\rm IGM}\sim0.1$ within a few kpc of the galaxy; ionization and
  heating rates fall as $r^{-2}$ (\S5.1, \S6 point v).
\item[\textbf{R4}] Extrapolated to a homogeneous population, the mean ionization
  rate is ${\rm d}n_{\rm ion}/{\rm d}t\sim8\times10^{-18}\,{\rm m^{-3}\,s^{-1}}$,
  ``of the order of the typical estimates of Leite et al.\ (2017) for CRs from
  SNe'' ($2\times10^{-18}\,{\rm m^{-3}\,s^{-1}}$) but far below Pop II UV
  ($\sim10^{-14}$) (\S6).
\item[\textbf{R5}] The heating rate is
  ${\rm d}{\rm Heat}/{\rm d}t\sim8\times10^{-20}\,{\rm keV\,m^{-3}\,s^{-1}}$,
  corresponding to a global $\Delta T\sim50$~K at $z\sim10$ --- of the same
  order as Leite et al.\ (2017) for CRs and Madau \& Fragos (2017) for X-rays
  (\S6).
\item[\textbf{R6}] Soft spectra are better ionizers and heaters; rates are
  insensitive to the high-energy cut-off and sensitive to the low-energy one
  (\S5.1, \S6 point vi).
\item[\textbf{R7}] Only $\sim5$ per cent of the source power is deposited in the
  IGM inside 1~Mpc; $\sim60$ per cent leaves the 1~Mpc volume (\S5.3).
\end{description}

% =====================================================================
\section{Catalogue of conflicting papers (folder \texttt{against/})}
% =====================================================================

Ordered by conflict strength, then by recency.

% ---------------------------------------------------------------------
\subsection{Dhandha, Fialkov, Gessey-Jones, et al.\ (2025) --- \HYP{} --- \textsc{direct} (strongest entry)}
% ---------------------------------------------------------------------
\noindent\textbf{Dhandha, J., Fialkov, A., Gessey-Jones, T., et al.\ 2025},
\emph{Narrowing the discovery space of the cosmological 21-cm signal using
multi-wavelength constraints}, MNRAS \textbf{544}, 1608. DOI:
\href{https://doi.org/10.1093/mnras/staf1736}{10.1093/mnras/staf1736};
arXiv:2508.13761; bibcode \texttt{2025MNRAS.544.1608D}.
\emph{Found via ADS full-text search; absent from the arXiv-only sweep.}\\
\textbf{File:} \texttt{against/Dhandha (2025) - Narrowing the discovery space\dots\ [arXiv 2508.13761].pdf}

\paragraph{Conflicts with:} \textbf{H4} (the warm, $T_{\rm IGM}\sim10^{4}$~K IGM
at $z=10$), and through it \textbf{R3}. This is the same conflict as the HERA
entry below, but stated \emph{at Douna et al.'s exact fiducial redshift} and
with a tighter bound.

\paragraph{Verbatim excerpt (Abstract)}
\begin{quote}\small
``From CXB and HERA data, we infer the IGM kinetic temperature to be
$T_{\rm K}(z = 15) \lesssim 7.7$~K, $2.5\,{\rm K} \lesssim T_{\rm K}(z = 10)
\lesssim 66\,{\rm K}$, and $20\,{\rm K} \lesssim T_{\rm K}(z = 6) \lesssim
2078\,{\rm K}$ at 95\% credible interval (CI).''
\end{quote}

\paragraph{Analysis.}
Douna et al.\ set $z=10$ as their working redshift (\S2.1) and evaluate the
recombination rate their electrons must outcompete at $T_{\rm IGM}\sim10^{4}$~K
(\S5.1). Dhandha et al.\ infer $T_{\rm K}(z=10)\lesssim66$~K at 95\% credibility
by combining the cosmic X-ray background, HERA power-spectrum upper limits, the
SARAS~3 non-detection and HST/JWST UV luminosity functions. That is a factor
$\sim150$ below the assumed value, and it brackets $z=10$ directly rather than
by interpolation. It sits much closer to Douna et al.'s own ``extremely cold''
comparison case ($T\sim30$~K, the CMB temperature at $z=10$), for which they
state the total recombination coefficient is ``around 100 times higher than the
case B recombination coefficient''. On the observationally favoured branch,
therefore, the recombination rate their electrons compete against is roughly two
orders of magnitude larger than in the scenario from which R3
($f_{\rm ion}\sim0.1$ within a few kpc) is quoted.

\paragraph{Contradictory data / ambiguity.} \flag{explicit}
Four qualifications, all material.
(i)~\textbf{This is an inference, not a measurement.} It is a posterior from a
Pop~II galaxy model conditioned on several data sets, and the authors close
their own abstract with the caveat: ``the results depend on our assumption of a
redshift-independent X-ray and radio efficiency of galaxies and on the exclusion
of a flexible model for Population III stars.'' Since Douna et al.'s sources are
\emph{explicitly} Pop~III-era microquasars, the excluded model class is the one
her scenario inhabits. This weakens the entry considerably and must be quoted
alongside the bound.
(ii)~It leans on the SARAS~3 non-detection (\S3.13), itself contested.
(iii)~It cites Douna et al.\ (2015) --- the metallicity paper --- \emph{in
support} of the strong metallicity dependence, i.e.\ it corroborates \textbf{H1}
while contradicting \textbf{H4}.
(iv)~$T_{\rm K}$ is a volume-averaged quantity. Douna et al.'s R3 is a statement
about gas within a few kpc of a source, which a global mean does not directly
constrain. \textbf{I did not recompute the local-versus-global correction and
this remains the weakest link in the argument.}

% ---------------------------------------------------------------------
\subsection{HERA Collaboration (2023) --- \RES{} and \HYP{} --- \textsc{direct}}
% ---------------------------------------------------------------------
\noindent\textbf{Abdurashidova, Z., Adams, T., Aguirre, J.~E., et al.\
(HERA Collaboration), 2023}, \emph{Improved Constraints on the 21\,cm EoR Power
Spectrum and the X-Ray Heating of the IGM with HERA Phase~I Observations},
ApJ \textbf{945}, 124. DOI:
\href{https://doi.org/10.3847/1538-4357/acaf50}{10.3847/1538-4357/acaf50};
arXiv:2210.04912.\\
\textbf{File:} \texttt{against/HERA Collaboration (2023) - Improved Constraints\dots\ [arXiv 2210.04912].pdf}

\paragraph{Conflicts with:} \textbf{H4} (the warm, $T_{\rm IGM}\sim10^{4}$~K
IGM at $z=10$) and, through it, \textbf{R3} (which depends on $\alpha_{\rm B}$
evaluated at that temperature).

\paragraph{Verbatim excerpt (\S7, ``Constraints on the IGM'')}
\begin{quote}\small
``The new HPD 95\% credible intervals on the spin and kinetic temperatures are
$4.7\,{\rm K} < T_S < 171.2\,{\rm K}$ and $3.2\,{\rm K} < T_K < 313.2\,{\rm K}$
at $z = 10.4$ and $15.6\,{\rm K} < T_S < 656.7\,{\rm K}$ and
$13.0\,{\rm K} < T_K < 4768\,{\rm K}$ at $z = 7.9$.''
\end{quote}

\paragraph{Second excerpt (Abstract)}
\begin{quote}\small
``we find that the intergalactic medium must have been heated above the
adiabatic cooling limit at least as early as $z = 10.4$, ruling out a broad set
of so-called `cold reionization' scenarios. If this heating is due to high-mass
X-ray binaries during the cosmic dawn, as is generally believed, our result's
99\% credible interval excludes the local relationship between soft X-ray
luminosity and star formation''
\end{quote}

\paragraph{Analysis.}
Douna et al.\ evaluate the recombination rate that their electrons must
outcompete at $T_{\rm IGM}\sim10^{4}$~K (\S5.1). HERA's 95\% upper bound at
$z=10.4$ --- the redshift closest to Douna's fiducial $z=10$, and slightly
\emph{later}, hence if anything warmer --- is $T_K < 313.2$~K, i.e.\ lower than
the assumed value by a factor $\gtrsim30$. Douna et al.\ themselves state the
consequence: ``A colder medium would imply a higher recombination rate, and thus
a lower reionization power of electrons'' (\S5.1). Their own ``cold'' bracket
uses $T\sim30$~K, where the recombination coefficient is ``around 100 times
higher than the case B recombination coefficient''. The observationally
permitted temperature at $z\approx10$ therefore sits in the lower part of their
bracketed range, not at the warm end on which R3 ($f_{\rm ion}\sim0.1$ within a
few kpc) is quoted. This entry is retained alongside Dhandha et al.\ (\S3.1)
because it is the primary observational result: Dhandha et al.\ tighten the
bound and place it at $z=10$ exactly, but only through an inference layer that
HERA's own limit does not require.

\paragraph{Contradictory data / ambiguity.} \flag{explicit}
This same paper cuts the other way on \textbf{H2}. HERA \emph{requires} heating
above adiabatic and excludes the local $L_{\rm X}$--SFR relation, which is an
argument that early accreting binaries were \emph{more}, not less, effective ---
supporting Douna's H1 while undermining her H2 (that X-rays are ``marginal at
most''). The paper also states of the weak-X-ray exclusion that ``that statement
is somewhat model dependent''. So: HERA contradicts H4 directly and H2
directly, but in doing so it strengthens H1. This entry must not be cited as a
blanket refutation.

% ---------------------------------------------------------------------
\subsection{Leite, Evoli, D'Angelo, Ciardi, Sigl \& Ferrara (2017) --- \RES{} --- \textsc{direct}}
% ---------------------------------------------------------------------
\noindent MNRAS \textbf{469}, 416--424. DOI:
\href{https://doi.org/10.1093/mnras/stx805}{10.1093/mnras/stx805};
arXiv:1703.09337.\\
\textbf{File:} \texttt{against/Leite (2017) - Do Cosmic Rays Heat the Early Intergalactic Medium [arXiv 1703.09337].pdf}

\paragraph{Conflicts with:} \textbf{R3} and \textbf{R4} (the ionization claims).

\paragraph{Verbatim excerpt (Abstract)}
\begin{quote}\small
``we show that CRs contribute negligibly to IGM ionization, but heat it
substantially, raising its temperature by $\Delta T = 10 - 200$~K by $z = 10$,
depending on the CR injection spectrum.''
\end{quote}

\paragraph{Second excerpt (\S4.1, ``Impact of CRs on reionization'')}
\begin{quote}\small
``The CR ionization rate is several orders of magnitude smaller than the UV
photoionization rate. This justifies the fact that CR were not included when we
described our reionization model in \S2.2.''
\end{quote}

\paragraph{Analysis.}
This is the single sharpest conflict in the audit, and it is internal to
Douna et al.'s own benchmarking. In \S6 they write that their extrapolated
ionization rate is ``of the order of the typical estimates of Leite et al.\
(2017) for CRs from SNe'' --- and then conclude (\S6, point v) that MQs ``would
be able to maintain by themselves large ionization fractions (of the order of
0.1) near the galaxies''. Leite et al., for a rate of the same order, conclude
the opposite: that the CR contribution to ionization is negligible and can be
dropped from the reionization model entirely. The two papers agree on the
number and disagree on what it means.

\paragraph{Contradictory data / ambiguity.} \flag{explicit}
Three qualifications. (i)~\textbf{Chronology}: Leite et al.\ (2017) precedes
Douna et al.\ (2018) and is cited by them; this is a tension Douna et al.\ were
aware of and did not resolve, not a subsequent refutation.
(ii)~\textbf{Carrier}: Leite et al.\ model CR \emph{protons} from Pop~II
supernovae; Douna et al.\ model CR \emph{electrons} from MQ jets. The rates are
compared by the original authors, so the comparison is theirs, but the physical
systems differ.
(iii)~\textbf{Scale}: Leite et al.\ compare against the \emph{global} UV
photoionization rate, whereas Douna et al.'s $f_{\rm ion}\sim0.1$ claim is
explicitly local (within a few kpc of the source). The two statements are not
strictly mutually exclusive; the conflict is in the interpretation of
significance, not in arithmetic. I could not resolve this cleanly and flag it
as unresolved.

% ---------------------------------------------------------------------
\subsection{Gessey-Jones, Fialkov, de Lera Acedo, Handley \& Barkana (2023) --- \RES{} --- \textsc{direct}}
% ---------------------------------------------------------------------
\noindent MNRAS \textbf{526}, 4262--4284. DOI:
\href{https://doi.org/10.1093/mnras/stad3014}{10.1093/mnras/stad3014};
arXiv:2304.07201.\\
\textbf{File:} \texttt{against/Gessey-Jones (2023) - Signatures of Cosmic Ray Heating in 21-cm Observables [arXiv 2304.07201].pdf}

\paragraph{Conflicts with:} \textbf{R4}, \textbf{R5} and \textbf{R7} (the
extrapolation from kpc-scale deposition to a homogeneous, cosmologically
relevant mean rate).

\paragraph{Verbatim excerpt (\S5, Discussion and Conclusions)}
\begin{quote}\small
``We find the spatial range of cosmic ray heating to be much shorter than that
of X-ray heating, even when cosmic rays are assumed to travel along straight
lines. This short-range nature manifests in the IGM heating being clustered
around regions of efficient star formation and results in a sharp contrast in
the 21-cm signal between localized heated regions of emission and vast unheated
regions of absorption.''
\end{quote}

\paragraph{Analysis.}
The clause ``even when cosmic rays are assumed to travel along straight lines''
is decisive, because rectilinear propagation is precisely Douna et al.'s
transport model (H6). Douna et al.\ nonetheless convert a kpc-scale, $r^{-2}$
deposition profile into a single homogeneous mean rate (R4, R5) by imagining
``a set of homogeneously distributed galaxies containing MQs''. Gessey-Jones
et al.\ show that this averaging step is exactly what CR heating does not
permit: the medium separates into hot islands and vast cold voids, so a
volume-averaged $\Delta T\sim50$~K does not describe the thermal state of the
IGM anywhere.

\paragraph{Contradictory data / ambiguity.} \flag{explicit}
Two, both material.
(i)~The same paper reports that ``cosmic rays with much higher energies up to
200 MeV have a non-negligible contribution, heating up the IGM on large scales
out to $\sim100$~cMpc away from their sources'' --- so the short-range
conclusion is spectrum-dependent, and Douna et al.'s hard tail (up to 1~PeV)
is not covered by it.
(ii)~Gessey-Jones et al.\ model \textbf{CR protons only}; they explicitly set
CR electrons aside (footnote~1: ``Cosmic ray electrons are however important to
consider when modelling SNR-produced radio backgrounds''). Their transport
conclusion is therefore not a computation about Douna et al.'s carrier. Grade
\textsc{direct} is assigned because the geometric argument is carrier-blind,
but this is a judgement call, not a demonstration.

% ---------------------------------------------------------------------
\subsection{Owen, Jacobsen, Wu \& Surajbali (2018) --- \HYP{} --- \textsc{direct}}
% ---------------------------------------------------------------------
\noindent MNRAS \textbf{481}, 666--687. DOI:
\href{https://doi.org/10.1093/mnras/sty2279}{10.1093/mnras/sty2279};
arXiv:1808.07837.\\
\textbf{File:} \texttt{against/Owen (2018) - Interactions between Ultra-High-Energy Particles and Protogalactic Environments [arXiv 1808.07837].pdf}

\paragraph{Conflicts with:} \textbf{H6} (omission of magnetic fields) and, as a
consequence, \textbf{R1}, \textbf{R2} and \textbf{R7} (escape fractions and the
energy budget leaving the galaxy).

\paragraph{Verbatim excerpt (Abstract)}
\begin{quote}\small
``We have found that protogalactic magnetic fields and their evolution can
significantly affect the energy transport and energy deposition processes of
cosmic rays. Within a Myr after the onset of star-formation the magnetic field
in a protogalaxy could attain a strength sufficient to confine all but the
highest energy particles within the galaxy.''
\end{quote}

\paragraph{Analysis.}
Douna et al.\ defer magnetic fields to a companion paper and compute
rectilinear escape from a 1~kpc galaxy. If a protogalaxy magnetizes within
$\sim1$~Myr to the point of confining all but the highest-energy particles,
then the escape fractions of R1--R2 --- and with them the $\sim65$ per cent of
injected power that Douna et al.\ report as reaching the IGM (R7) --- are
computed in a regime that does not obtain. Note the direction: confinement
would suppress the escape of exactly the intermediate-energy electrons that
Douna et al.\ identify as the reservoir feeding the low-energy ionizing
population (R2, R6).

\paragraph{Contradictory data / ambiguity.} \flag{explicit}
The same abstract preserves a window for Douna et al.'s scenario: ``there is a
window before the protogalaxy is fully magnetised, energetic particles could
stream across the galaxy freely, delivering energy into the circumgalactic and
intergalactic medium.'' Whether MQ jets operate inside or outside that window is
not determined by either paper. Owen et al.\ also model \emph{hadronic}
particles; the confinement argument is magnetic-rigidity-based and so applies to
electrons of the same rigidity, but the paper performs no electron calculation.

% ---------------------------------------------------------------------
\subsection{Kaur, Qin \& Mesinger (2022) --- \HYP{} --- \textsc{moderate}}
% ---------------------------------------------------------------------
\noindent MNRAS \textbf{513}, 5097. DOI:
\href{https://doi.org/10.1093/mnras/stac1226}{10.1093/mnras/stac1226};
arXiv:2203.10851; bibcode \texttt{2022MNRAS.513.5097K}.
\emph{Found via ADS topical search.}\\
\textbf{File:} \texttt{against/Kaur (2022) - The 21-cm signal from the cosmic dawn - metallicity dependence\dots\ [arXiv 2203.10851].pdf}

\paragraph{Conflicts with:} \textbf{H1} --- not its direction, but the precision
with which Douna et al.\ deploy it.

\paragraph{Verbatim excerpt (Abstract)}
\begin{quote}\small
``For our fiducial models, galaxies with star formation rates of order
$10^{-3}$--$10^{-1}\,M_\odot$~yr$^{-1}$ and metallicities of order
$10^{-3}$--$10^{-2}\,Z_\odot$ are the dominant contributors to the X-ray
background (XRB) during this period. Different $L_{\rm X}$/SFR--$Z$ relations
result in factors of $\sim3$ differences in these ranges, as well as in the mean
IGM temperature and the large-scale 21-cm power, at a given redshift.''
\end{quote}

\paragraph{Analysis.}
Douna et al.\ invoke the metallicity enhancement (H1) qualitatively and then
adopt a single fiducial $L_{\rm k}=10^{40}\,{\rm erg\,s^{-1}}$ with no
uncertainty band. Kaur et al.\ show that merely \emph{choosing between published
$L_{\rm X}$/SFR--$Z$ relations} moves the mean IGM temperature by a factor
$\sim3$. Since all of Douna et al.'s rates scale linearly with $L_{\rm k}$
(\S5), an unquantified factor of a few propagates directly into R3--R5. Together
with Sartorio et al.\ (\S3.10), which adds an IMF-driven spread of orders of
magnitude, this is the case for attaching an error bar to the fiducial
luminosity rather than quoting single-valued rates.

\paragraph{Contradictory data / ambiguity.} \flag{explicit}
Kaur et al.\ model the \emph{X-ray} output of HMXBs, not jet kinetic power; the
factor-3 spread is not automatically the spread in $L_{\rm k}$. The paper
affirms the metallicity dependence itself, so this contradicts the treatment of
H1 as a precise input, not H1 as a claim.

% ---------------------------------------------------------------------
\subsection{Jana \& Nath (2018) --- \RES{} --- \textsc{moderate}}
% ---------------------------------------------------------------------
\noindent MNRAS \textbf{479}, 153--161. DOI:
\href{https://doi.org/10.1093/mnras/sty1481}{10.1093/mnras/sty1481};
arXiv:1806.01295.\\
\textbf{File:} \texttt{against/Jana (2018) - Cosmic ray heating of intergalactic medium - patchy or uniform [arXiv 1806.01295].pdf}

\paragraph{Conflicts with:} \textbf{R3} (the localized, kpc-scale picture) and
\textbf{H6}/\textbf{H8} (rectilinear transport through a neutral IGM).

\paragraph{Verbatim excerpt (\S5, Summary)}
\begin{quote}\small
``In the case of CRs from mini-haloes at $z\sim10$--20, we put an upper bound on
the diffusion coefficient ($D\le1\times10^{26}$~cm$^{2}$~s$^{-1}$ for $z\sim10$
and $D\le5$--$6\times10^{26}$~cm$^{2}$~s$^{-1}$ for $z\sim20$) for which the
heating is inhomogeneous, after comparing our temperature profiles with the
estimate of global temperature increase. Given the expected scaling of the
diffusion coefficient with redshift, this bound suggests uniform heating, both
at $z\sim10$ and 20.''
\end{quote}
\flag{transcription} The extractor rendered the printed ``$\le$'' as the digit
``6''; restored above. Exponents were reconstructed from the printed form
(``$1\times10^{26}$~cm$^{2}$~s$^{-1}$'').

\paragraph{Second excerpt (\S5, Summary)}
\begin{quote}\small
``The surrounding medium of galaxies at high redshift is likely to be ionized on
account of star formation in the galaxy. Therefore, the heating by CRs will
proceed in a different manner than previously considered case of neutral IGM
heating.''
\end{quote}

\paragraph{Analysis.}
Douna et al.'s central spatial claim (R3) is that the interesting action is
\emph{local} --- $f_{\rm ion}\sim0.1$ within a few kpc, falling as $r^{-2}$.
Jana \& Nath, using the diffusion coefficient expected from the scaling
$D\propto B^{-1/3}$ with a $\sim10^{-9}$~G IGM field, conclude that CR heating
at $z\sim10$ should be \emph{uniform}, not patchy. Their second excerpt also
conflicts with the framing of H8: they argue the near-galaxy medium is already
photoionized, so the neutral-IGM energy-partition calculation Douna et al.\
perform is not the applicable one.

\paragraph{Contradictory data / ambiguity.} \flag{explicit}
The authors immediately disclaim their own conclusion: ``But the uncertainties
in CR parameters (spectrum, lower energy limit of emerging CRs, diffusion
coefficient and its dependence on magnetic field) and magnetic field at high
redshift precludes any firm conclusion.'' Note also that Jana \& Nath and
Gessey-Jones et al.\ (previous entry) reach \emph{opposite} conclusions on
patchy-vs-uniform. This is an open disagreement in the literature, not a settled
result against Douna et al.; what both share is that neither supports her
combination of local structure \emph{and} a meaningful global average.
Jana \& Nath's first summary bullet --- ``It is not easy for low energy CRs to
escape relatively massive galaxy'' --- \emph{agrees} with Douna et al.'s R1.

% ---------------------------------------------------------------------
\subsection{Yokoyama \& Ohira (2023) --- \HYP{} (model) --- \textsc{moderate}}
% ---------------------------------------------------------------------
\noindent MNRAS \textbf{523}, 3671--3677. DOI:
\href{https://doi.org/10.1093/mnras/stad1596}{10.1093/mnras/stad1596};
arXiv:2302.07028.\\
\textbf{File:} \texttt{against/Yokoyama (2023) - Resistive Heating Induced by Streaming Cosmic Rays Around a Galaxy in the Early Universe [arXiv 2302.07028].pdf}

\paragraph{Conflicts with:} \textbf{H6} and the process inventory of \textbf{\S2.1},
and hence with the mechanism attribution behind \textbf{R5}.

\paragraph{Verbatim excerpt (\S4, Summary)}
\begin{quote}\small
``We revisited the influence of streaming CRs on the heating of the IGM. We
clarified that resistive heating induced by CRs becomes dominant over the other
mechanisms in the vicinity of a galaxy at $r\lesssim10^{2}$~kpc in early phases.
The temperature reaches $10^{4}$~K at $r = 1$~kpc in $2\times10^{6}$~yr.''
\end{quote}

\paragraph{Analysis.}
Douna et al.\ enumerate their heating channels explicitly in \S2.1:
collisional ionization and excitation, elastic $e^-e^-$ scattering,
bremsstrahlung, inverse Compton and direct Compton. Yokoyama \& Ohira identify a
channel absent from that list --- Ohmic dissipation of the thermal return
current driven by the streaming CR beam --- and find it dominant over the
collisional channels precisely in the spatial range ($r\lesssim100$~kpc) where
Douna et al.\ compute their profiles. The attribution of the heating in R5 to
$e^-e^-$ scattering is therefore incomplete for a streaming population.

\paragraph{Contradictory data / ambiguity.} \flag{explicit --- direction matters}
This conflict does \emph{not} weaken Douna et al.'s conclusion; it points the
same way with a different mechanism, and would if anything raise the heating
rate. It is included because the user's criterion covers conflicts with the
model as well as with the numbers, and because it falsifies Douna et al.'s
claim (H7, \S6) that their omissions are uniformly ``conservative'' in the sense
of being physically complete-but-underestimated: a missing dominant channel is
a different kind of gap. Note also that resistive heating requires a streaming
beam and a return current, i.e.\ a magnetized medium --- the very ingredient
Douna et al.\ set aside. Yokoyama \& Ohira model \textbf{CR protons}.

% ---------------------------------------------------------------------
\subsection{Owen, Jin, Wu \& Ferreras (2019) --- \RES{} --- \textsc{moderate}}
% ---------------------------------------------------------------------
\noindent MNRAS \textbf{484}, 1645--1671. DOI:
\href{https://doi.org/10.1093/mnras/stz060}{10.1093/mnras/stz060};
arXiv:1901.01411.\\
\textbf{File:} \texttt{against/Owen (2019) - Hadronic Interactions of Energetic Charged Particles in Protogalactic Outflow Environments [arXiv 1901.01411].pdf}

\paragraph{Conflicts with:} \textbf{R4}, \textbf{R5}, \textbf{R7} (the spatial
extent of energy deposition).

\paragraph{Verbatim excerpt (Abstract)}
\begin{quote}\small
``In a purely diffusive transport scenario, we find the peak heating rate
reaches $10^{-26}$~erg~cm$^{-3}$~s$^{-1}$ at the base of the outflow where the
wind is driven by core-collapse supernovae at an event rate of
0.1~yr$^{-1}$, but does not extend beyond 2~kpc. In the advection limit, the
peak heating rate is reduced to $10^{-28}$~erg~cm$^{-3}$~s$^{-1}$, but its
extent can reach to tens of kpc.''
\end{quote}

\paragraph{Analysis.}
Douna et al.\ follow electrons to 1~Mpc and report that $\sim60$ per cent of the
emitted power still leaves that volume (R7). Owen et al.\ find that, once
realistic diffusive or advective transport through a protogalactic outflow is
imposed, CR energy deposition terminates at 2~kpc (diffusive) or tens of kpc
(advective) --- one and a half to three orders of magnitude short of the volume
over which Douna et al.\ compute and then extrapolate.

\paragraph{Contradictory data / ambiguity.} \flag{explicit}
The comparison is not like-for-like in two respects: Owen et al.\ treat
\emph{hadronic} CRs in a supernova-driven \emph{outflow}, whereas Douna et al.\
treat electrons injected by a jet into a static homogeneous medium; and the two
papers use different energy-loss physics (pion production vs.\ Coulomb and
inverse Compton). The conflict is over what transport regime is appropriate, and
neither paper adjudicates it. \textbf{This entry was not independently
recomputed.}

% ---------------------------------------------------------------------
\subsection{Sartorio, Fialkov, Hartwig, et al.\ (2023) --- \HYP{} --- \textsc{moderate}}
% ---------------------------------------------------------------------
\noindent MNRAS \textbf{521}, 4039--4055. DOI:
\href{https://doi.org/10.1093/mnras/stad697}{10.1093/mnras/stad697};
arXiv:2303.03435.\\
\textbf{File:} \texttt{against/Sartorio (2023) - Population III X-ray Binaries and their Impact on the Early Universe [arXiv 2303.03435].pdf}

\paragraph{Conflicts with:} \textbf{H1} and \textbf{H5} (the assumption that a
single fiducial $L_{\rm k}=10^{40}\,{\rm erg\,s^{-1}}$ characterizes the
high-$z$ MQ population).

\paragraph{Verbatim excerpt (Abstract)}
\begin{quote}\small
``Specifically, we find that the uniformity and strength of the X-ray feedback
in the intergalactic medium is strongly dependant on the IMF. Bottom-heavy IMFs
result in a smoother distribution of XRBs, but have a luminosity orders of
magnitude lower than more top-heavy IMFs. Top-heavy IMFs lead to more spatially
uneven, albeit strong, X-ray emission.''
\end{quote}

\paragraph{Second excerpt (\S5.5, Halo Absorption)}
\begin{quote}\small
``soft photons with energies below 0.1 keV are completely absorbed and, thus,
contribute nothing to IGM heating and ionization.''
\end{quote}

\paragraph{Analysis.}
Douna et al.\ fix $L_{\rm k}=10^{40}\,{\rm erg\,s^{-1}}$ and note that ``all the
resulting spectra are linear in this parameter'' (\S5). Sartorio et al.\ show
that the population-integrated luminosity of primordial XRBs spans
\emph{orders of magnitude} depending on the Pop~III IMF, which is unconstrained.
Because Douna et al.'s results scale linearly with $L_{\rm k}$, the
uncertainty in H5 propagates linearly into R3--R5, an uncertainty band Douna
et al.\ do not carry. This is a conflict with the premise's precision, not with
its sign.

\paragraph{Contradictory data / ambiguity.} \flag{explicit}
Sartorio et al.\ treat the \emph{X-ray} output of Pop~III XRBs, not jet kinetic
power; the mapping from one to the other is model-dependent and is not made in
either paper. Their second excerpt is a conflict with \textbf{H2} in the
opposite direction from HERA's: it supports Douna et al.'s premise that the soft
X-ray channel is ineffective. \textbf{The two excerpts from this single paper
therefore cut in opposite directions} and the entry should be cited with that
stated.

% ---------------------------------------------------------------------
\subsection{Lehmer, Eufrasio, Basu-Zych, et al.\ (2021) --- \HYP{} --- \textsc{moderate}}
% ---------------------------------------------------------------------
\noindent ApJ \textbf{907}, 17. DOI:
\href{https://doi.org/10.3847/1538-4357/abcec1}{10.3847/1538-4357/abcec1};
arXiv:2011.09476.\\
\textbf{File:} \texttt{against/Lehmer (2021) - The Metallicity Dependence of the HMXB Luminosity Function [arXiv 2011.09476].pdf}

\paragraph{Conflicts with:} \textbf{H1} (the metallicity premise, as stated
without luminosity qualification).

\paragraph{Verbatim excerpt (Abstract)}
\begin{quote}\small
``We demonstrate that the SFR-normalized HMXB XLF shows clear trends with
metallicity, with steadily increasing numbers of luminous and ultraluminous
X-ray sources ($\log L\,({\rm erg\,s^{-1}}) = 38$--40.5) with declining
metallicity. However, we find that the low-luminosity
($\log L\,({\rm erg\,s^{-1}}) = 36$--38) HMXB XLF appears to show a nearly
constant SFR scaling and slope with metallicity.''
\end{quote}

\paragraph{Analysis.}
Douna et al.\ state H1 without restriction: ``there is strong evidence that
supports that XRBs are more numerous in low-metallicity environments'' (\S1),
citing seven references including Douna et al.\ (2015). Lehmer et al.\ show the
enhancement is confined to $L_{\rm X}\gtrsim10^{38}\,{\rm erg\,s^{-1}}$ and is
absent below it. The premise survives for the ULX-class systems relevant to
$L_{\rm k}=10^{40}\,{\rm erg\,s^{-1}}$, but not as the general population
statement Douna et al.\ make.

\paragraph{Contradictory data / ambiguity.} \flag{explicit}
This is a refinement that mostly \emph{supports} Douna et al.'s specific use
case, since her fiducial source is at the luminous end. Graded
\textsc{moderate}, not \textsc{direct}, for that reason. Note that Lehmer et
al.\ (2021 and 2024) cite Douna et al.\ (2015) --- the metallicity paper --- as
supporting evidence, not the 2018 paper under audit.

% ---------------------------------------------------------------------
\subsection{Sotomayor Checa, Romero \& Bosch-Ramon (2021) --- \HYP{} --- \textsc{moderate}}
% ---------------------------------------------------------------------
\noindent Ap\&SS \textbf{366}, 1. DOI:
\href{https://doi.org/10.1007/s10509-020-03911-5}{10.1007/s10509-020-03911-5};
arXiv:2012.11664 (posted 2020).\\
\textbf{File:} \texttt{against/Sotomayor Checa (2020) - Equatorial outflows driven by jets in Population III microquasars [arXiv 2012.11664].pdf}

\paragraph{Conflicts with:} \textbf{H5} / \textbf{H3} (that a MQ jet delivers a
clean power-law electron population into the ISM).

\paragraph{Verbatim excerpt (\S5, conclusions)}
\begin{quote}\small
``The geometry we have adopted for the wind of the accretion disk is a purely
axisymmetrical flow. In this way, we can analyze more adequately the production
of a jet-induced equatorial wind. However, a more realistic geometry should take
into account a wind that interacts collisionally with the jet. This would lead to
shocks in the jet and to a decrease in jet velocity, likely reducing the capacity
of the jet to emerge from the binary system.''
\end{quote}

\paragraph{Second excerpt (\S5, conclusions)}
\begin{quote}\small
``When $L_{\rm jet} \ll L_{\rm wind}$ (Case 3) or $L_{\rm jet}\sim L_{\rm wind}$
(Case 2), jets keep narrow while the winds do not get significantly deflected
sidewards. In such a case, the system is a powerful source of UV radiation and
soft X-rays emitted by the wind''
\end{quote}

\paragraph{Analysis.}
Relativistic hydrodynamic simulations by the group that built the Pop~III MQ
model show that in two of three power regimes the system's output is dominated
by wind UV and soft X-rays rather than by jet particles, and that a more
realistic jet--wind geometry would likely impede the jet from leaving the binary
at all. Douna et al.\ take the injection of jet electrons into the ISM as given
(\S2.1, ``We assume that electrons are produced in the jets of MQs and injected
into the ISM, as discussed by Heinz \& Sunyaev 2002'').

\paragraph{Contradictory data / ambiguity.} \flag{explicit --- weak entry}
The conflicting statement is the authors' \emph{own caveat on their own model},
phrased with ``likely'', not a demonstrated result; their abstract concludes the
opposite overall (``Population III microquasars might inject gamma rays and
relativistic particles into the early intergalactic medium, contributing to its
reionization''). This entry is the weakest in the catalogue and should be cited
as a modelling caution, not as a contradiction.

% ---------------------------------------------------------------------
\subsection{Yang, Long \& Hirata (2024) --- \RES{} (relevance) --- \textsc{indirect}}
% ---------------------------------------------------------------------
\noindent ApJ \textbf{965}, 111. DOI:
\href{https://doi.org/10.3847/1538-4357/ad3237}{10.3847/1538-4357/ad3237};
arXiv:2311.18721 (posted 2023).\\
\textbf{File:} \texttt{against/Yang (2023) - Compton scattering of electrons in the intergalactic medium [arXiv 2311.18721].pdf}

\paragraph{Conflicts with:} \textbf{H3} and the broader relevance claim behind
\textbf{R5} (that CR electrons are an energetically significant IGM component).

\paragraph{Verbatim excerpt (Abstract)}
\begin{quote}\small
``At mean density, the ratio of cosmic ray electron to thermal pressure in the
IGM $P_{\rm CRe}/P_{\rm th}$ is 0.3\% at $z=2$, rising to 1.0\% at $z=1$, and
1.8\% at $z=0.1$ (considering only the cosmic rays produced locally by Compton
scattering).''
\end{quote}

\paragraph{Analysis.}
The only paper found that computes the CR-electron energy budget of the IGM as
such. It finds the population sub-percent in pressure relative to the thermal
IGM and \emph{decreasing} toward high redshift (energy density peaks at
$z\approx2$).

\paragraph{Contradictory data / ambiguity.} \flag{explicit --- do not overclaim}
This is a weak entry and the mismatch is severe: the calculation runs
$z\le6$, not $z=10$; the CR electrons are produced by Compton scattering of
the extragalactic background, not injected by jets; and the parenthetical
``considering only the cosmic rays produced locally by Compton scattering''
explicitly excludes source-injected populations of the kind Douna et al.\ model.
It is included because it is the closest existing quantitative statement on the
relevance of IGM CR electrons, and because its redshift trend runs against the
assumption that this component grows in importance at cosmic dawn.
\textbf{No extrapolation of their numbers to $z=10$ was attempted and none
should be inferred.}

% ---------------------------------------------------------------------
\subsection{Bera, Samui \& Datta (2023) --- \HYP{} (attribution) --- \textsc{indirect}}
% ---------------------------------------------------------------------
\noindent MNRAS \textbf{519}, 4869--4883. DOI:
\href{https://doi.org/10.1093/mnras/stac3814}{10.1093/mnras/stac3814};
arXiv:2202.12308 (posted 2022).\\
\textbf{File:} \texttt{against/Bera (2022) - Impact of cosmic rays on the global 21-cm signal during cosmic dawn [arXiv 2202.12308].pdf}

\paragraph{Conflicts with:} \textbf{H2}/\textbf{H3} (the attribution of the
required CR heating to MQ electrons rather than SN protons) and the process
inventory of \textbf{\S2.1}.

\paragraph{Verbatim excerpt (Abstract)}
\begin{quote}\small
``The low energy cosmic ray protons from Pop III supernovae can escape from
minihalos and heat the IGM via collision and ionization of hydrogen. Furthermore,
high energy protons generated in Pop II supernovae can escape the hosting halos
and heat the IGM via magnetosonic Alfv\'en waves. We show that the heating due to
these cosmic ray particles can significantly impact the IGM temperature and hence
the global 21-cm signal at $z\sim14-18$.''
\end{quote}

\paragraph{Analysis.}
Two points of friction. First, the whole cosmic-dawn CR heating budget is
assigned to supernova protons, with no MQ contribution --- an attribution
conflict with Douna et al.'s framing that MQ electrons are the carrier ``expected
to make the largest contribution'' (H3). Second, the Alfv\'en-wave
(magnetosonic) transfer channel is absent from Douna et al.'s process list
(\S2.1) and, like the resistive channel of Yokoyama \& Ohira, requires the
magnetic field she omits (H6).

\paragraph{Contradictory data / ambiguity.} \flag{explicit}
Bera et al.\ nowhere assert that MQs are unimportant; they simply do not model
them. Absence of a term in someone else's budget is weak evidence against it.
The entry is retained for the missing-channel point, which is concrete.

% ---------------------------------------------------------------------
\subsection{Singh, Nambissan T., Subrahmanyan, et al.\ (SARAS 3; 2022) --- \HYP{} (context) --- \textsc{indirect}}
% ---------------------------------------------------------------------
\noindent Nature Astronomy \textbf{6}, 607--617. DOI:
\href{https://doi.org/10.1038/s41550-022-01610-5}{10.1038/s41550-022-01610-5};
arXiv:2112.06778 (posted 2021).\\
\textbf{File:} \texttt{against/Singh (2022) - SARAS 3 - On the detection of a cosmic dawn signal in the radio background [arXiv 2112.06778].pdf}

\paragraph{Conflicts with:} the observational context invoked for
\textbf{H2}/\textbf{H4} --- specifically the EDGES-driven case that cosmic dawn
requires non-standard heating or radio backgrounds.

\paragraph{Verbatim excerpt (Abstract)}
\begin{quote}\small
``We report a radiometer measurement of the spectrum of the radio sky in the
55--85~MHz band, which shows that the profile found by Bowman et al.\ in data
taken with the Experiment to Detect the Global Epoch of Reionization Signature
(EDGES) low-band instrument is not of astrophysical origin; their best-fitting
profile is rejected with 95.3\% confidence.''
\end{quote}

\paragraph{Analysis.}
Douna et al.\ (2018) predate EDGES and do not invoke it, so this is
\emph{not} a conflict with their paper as written. It is included because the
downstream literature that carries her line of work forward --- Mirabel (2019),
Mirabel \& Rodr\'iguez (2022), both in the parent folder --- rests substantially
on the EDGES absorption being real. If the user's paper builds on that thread,
this is the constraint that has to be addressed.

\paragraph{Contradictory data / ambiguity.} \flag{explicit}
The EDGES/SARAS~3 disagreement is unresolved in the literature. A rejection at
95.3\% confidence is not a refutation, and the two experiments have not been
reconciled. \textbf{This entry is context, not evidence against Douna et al.}

% =====================================================================
\section{Conflicts found inside the parent folder}
\label{sec:infolder}
% =====================================================================

These four papers are already in the parent folder (they are among the 14 that
cite Douna et al.\ 2018) and were \emph{not} duplicated into \texttt{against/}.
Each contains a statement in conflict with a Douna claim while citing her
neutrally or approvingly elsewhere.

\begin{longtable}{@{}p{3.3cm}p{1.4cm}p{10.3cm}@{}}
\toprule
\textbf{Paper} & \textbf{Target} & \textbf{Verbatim conflicting excerpt} \\
\midrule
\endhead
Sazonov \& Khabibullin (2018), MNRAS \textbf{476}, 2530; arXiv:1712.04831
& \HYP{} \textbf{H2}
& ``This implies that ULXs (and hence HMXBs as a whole) can induce a stronger
heating of the IGM in the early Universe compared to estimates neglecting X-ray
feedback.'' (Summary) \\[4pt]
\multicolumn{3}{@{}p{15.2cm}@{}}{\footnotesize\emph{Analysis:} Douna et al.\
motivate the CR-electron channel by asserting the X-ray channel is ``marginal at
most''. Sazonov \& Khabibullin show the soft X-ray escape fraction from a
ULX-hosting minihalo rises from $\sim$20--50\% to $\sim$30--80\% over the ULX
lifetime, i.e.\ X-rays are \emph{more} effective than the estimates Douna et al.\
rely on. \flag{ambiguity} The same Summary limits this to
$M\lesssim3\times10^{8}\,M_{\odot}$ haloes and concludes the population-averaged
escape fraction ``probably remained low, $\lesssim30$\%, at these epochs''.} \\
\midrule
Bosch-Ramon (2018), A\&A \textbf{617}, L3; arXiv:1808.08911
& \HYP{} \textbf{H3}
& ``only the most energetic cosmic rays can diffuse out of the radio lobes;
otherwise, they stay in the lobes and cool through adiabatic losses. Secondly,
it is not known which fraction of the jet energy is in the form of very
energetic cosmic rays.'' (\S1) \\[4pt]
\multicolumn{3}{@{}p{15.2cm}@{}}{\footnotesize\emph{Analysis:} Written of AGN
lobes, but it is the reason Bosch-Ramon abandons escaping CRs and works with
\emph{in situ} lobe electrons radiating via inverse Compton instead. Douna et
al.\ assume escape. \flag{ambiguity} The scale argument does not transfer
cleanly from Mpc-scale AGN lobes to pc-scale MQ jets; Bosch-Ramon cites Douna et
al.\ approvingly two sentences earlier.} \\
\midrule
Escobar, Pellizza \& Romero (2022), A\&A \textbf{665}, A61; arXiv:2207.08633
& \HYP{} \textbf{H3}
& ``We focus on the secondary protons product of neutron beta decay. Electrons
and neutrinos are also emitted in the process but they carry a negligible
fraction of the progenitor neutron energy ($\sim0.1$\%).'' (\S2) \\[4pt]
\multicolumn{3}{@{}p{15.2cm}@{}}{\footnotesize\emph{Analysis:} The same group's
2021 paper (A\&A 650, A136) explicitly set out to ``support the claims of Douna
et al.\ (2018)''; by 2022 the electron component has been dropped as
energetically negligible in the neutron-escape channel, and the mature version
of the mechanism is proton-dominated. \flag{ambiguity} This concerns the
neutron-decay channel specifically, not shock-accelerated jet electrons, which
is what Douna et al.\ model.} \\
\midrule
Mirabel \& Rodr\'iguez (2022), New Astron.\ Rev.\ \textbf{94}, 101642; arXiv:2203.12741
& \HYP{} \textbf{H1}
& ``The unreasonable large number of sources of this type that are required to
match the onset of the EDGES trough would exceed the expected maximum total mass
of stars already formed at that epoch by factors of $\sim10^{6}$ for sources like
Cygnus X-1'' (\S X, conclusion 4) \\[4pt]
\multicolumn{3}{@{}p{15.2cm}@{}}{\footnotesize\emph{Analysis:} A co-author of
Douna et al.\ (2018) arguing that BH-HMXB microquasars of the known Galactic type
cannot, in the required numbers, drive the cosmic-dawn signal. \flag{ambiguity}
The argument is specifically about matching the EDGES trough onset via a radio
background --- a different observable from IGM heating --- and the paper's
alternative is IMBHs, not the abandonment of the microquasar picture.} \\
\bottomrule
\end{longtable}

% =====================================================================
\section{Screened, no conflict found}
\label{sec:screened}
% =====================================================================

Downloaded, read, and moved to \texttt{screened\_no\_conflict/} because
full-text inspection produced no statement in conflict with any claim
in Section~\ref{sec:claims}. Listed so the search is auditable.

\begin{itemize}[leftmargin=*,itemsep=3pt]
\item \textbf{Sazonov \& Khabibullin (2017)}, Astron.\ Lett.\ \textbf{43}, 211;
      arXiv:1612.01262 --- \emph{Preheating of the early universe by radiation
      from high-mass X-ray binaries}. Found by ADS full-text search and
      \textbf{initially catalogued as a conflict with H2 on the strength of its
      abstract}, which reads that HMXB X-rays ``could significantly heat
      \dots\ the intergalactic medium by $z\sim10$'' under stated conditions.
      Reading \S4 reverses that: ``Our conclusion that the Universe could not be
      significantly heated by X-rays from HMXBs before $z\sim10$ is similar to
      that reached by Madau \& Fragos (2016).'' The paper therefore
      \textbf{supports} Douna et al.'s premise H2 rather than contradicting it,
      and independently corroborates the Madau \& Fragos (2017) result she
      cites. \flag{This is a worked example of why abstract-level screening is
      not sufficient; the conditional in the abstract inverts in the
      conclusions.}
\item \textbf{Lehmer, Monson, Eufrasio, et al.\ (2024)}, arXiv:2410.19901 ---
      updated metallicity/SFH framework for XRB populations. Cites Douna et al.\
      (2015) as \emph{supporting} the $L_{\rm X}$/SFR--metallicity
      anti-correlation. No conflict; supersedes Lehmer et al.\ (2021) without
      reversing it.
\item \textbf{Leong \& Meiksin (2025)}, arXiv:2509.10255 --- 21-cm signature of
      X-ray heated haloes. Its ``most of the heating will be uniform on
      measurable scales'' reads as a conflict with Douna et al.'s localized
      picture, but the paper's substance (bright sources produce detectable warm
      rings) \emph{supports} the existence of a local heating zone. Net: no
      clear conflict. \flag{judgement call --- reconsider if the local/global
      distinction becomes central to your argument.}
\item \textbf{Dasgupta, Liu, Fialkov, et al.\ (2026)}, arXiv:2604.11542 ---
      stochastic Pop~III XRB luminosities in 21cmSPACE. Finds ``The impact of
      stochasticity on the global 21-cm signal and on the large-scale power
      spectrum is found to be negligible.'' Orthogonal to Douna et al.
\end{itemize}

\subsection*{Assessed from ADS abstracts, found supporting, not downloaded}

Two further papers surfaced by the ADS metallicity search were checked and found
to \emph{corroborate} H1, so no PDF was retrieved:

\begin{itemize}[leftmargin=*,itemsep=2pt]
\item \textbf{Fornasini, Kriek, Sanders, et al.\ (2019)}, ApJ \textbf{885}, 65
      (MOSDEF, $z\sim2$): ``Splitting our sample by $Z$, we find that
      $L_{\rm X}$/SFR and $Z$ are anticorrelated with 97\% confidence.''
      They add that ``the statistical significance of the observed trends is
      weak owing to the low X-ray statistics.''
\item \textbf{Fornasini, Civano \& Suh (2020)}, MNRAS \textbf{495}, 771:
      ``We find no significant variation of the $L_{\rm X}$--SFR--$Z$ relation
      with redshift. Our results provide further evidence that the $Z$
      dependence of HMXBs is responsible for the redshift evolution of
      $L_{\rm X}$/SFR.''
\end{itemize}
\flag{Assessed at abstract level only.} If H1 becomes load-bearing in your
paper, read both in full --- see the Sazonov \& Khabibullin (2017) reversal
above for why.

% =====================================================================
\section{Global assessment and confidence flags}
\label{sec:global}
% =====================================================================

\subsection{Summary table}

\begin{longtable}{@{}p{4.6cm}p{1.9cm}p{2.0cm}p{5.6cm}@{}}
\toprule
\textbf{Paper} & \textbf{Target} & \textbf{Strength} & \textbf{Claim(s) contradicted} \\
\midrule
\endhead
Dhandha et al.\ (2025)        & \HYP{}        & direct   & H4 ($T_K(z{=}10)\lesssim66$~K); R3 \\
HERA Collab.\ (2023)          & \HYP{}+\RES{} & direct   & H4 (warm IGM); H2 (opposite sense); R3 \\
Leite et al.\ (2017)          & \RES{}        & direct   & R3, R4 (ionization is negligible) \\
Gessey-Jones et al.\ (2023)   & \RES{}        & direct   & R4, R5, R7 (global averaging invalid) \\
Owen et al.\ (2018)           & \HYP{}        & direct   & H6 (no B field); R1, R2, R7 \\
Jana \& Nath (2018)           & \RES{}        & moderate & R3 (patchy); H6, H8 \\
Yokoyama \& Ohira (2023)      & \HYP{}        & moderate & \S2.1 process list; H6, H7 \\
Owen et al.\ (2019)           & \RES{}        & moderate & R4, R5, R7 (deposition range) \\
Sartorio et al.\ (2023)       & \HYP{}        & moderate & H5 ($L_{\rm k}$ uncertainty); H1 \\
Lehmer et al.\ (2021)         & \HYP{}        & moderate & H1 (only at $L>10^{38}$) \\
Kaur et al.\ (2022)           & \HYP{}        & moderate & H1 (factor $\sim3$ on $T_{\rm IGM}$) \\
Sotomayor Checa et al.\ (2021)& \HYP{}        & moderate & H3, H5 (jet may not emerge) \\
Yang et al.\ (2024)           & \RES{}        & indirect & H3 (CRe energetically minor) \\
Bera et al.\ (2023)           & \HYP{}        & indirect & H2, H3; missing Alfv\'en channel \\
Singh et al.\ (2022)          & \HYP{}        & indirect & EDGES context only \\
\bottomrule
\end{longtable}

\subsection{What the audit actually establishes}

\begin{enumerate}[leftmargin=*,itemsep=4pt]
\item \textbf{No published rebuttal of Douna et al.\ (2018) exists.}
      All 18 papers citing it --- the ADS list, which is a strict superset of
      what Semantic Scholar, OpenAlex, OpenCitations and INSPIRE-HEP return ---
      cite it neutrally or approvingly. No author has challenged its Monte
      Carlo, its escape fractions, or its rate estimates in print.
      \flag{confidence: high} --- this rests on the authoritative citation list
      and a full reading of the citation context in each.

\item \textbf{The strongest conflict is with a premise, not a calculation.}
      Dhandha et al.\ (2025) infer $T_K(z=10)\lesssim66$~K at 95\% credibility,
      and HERA (2023) measures $T_K < 313.2$~K at $z=10.4$ (95\%), against the
      assumed $T_{\rm IGM}\sim10^{4}$~K. Douna et al.\ bracketed this themselves
      and stated the direction of the effect, so the paper is not invalidated
      --- but the headline number R3 ($f_{\rm ion}\sim0.1$) is quoted from the
      warm branch, and the observationally permitted branch is the cold one.
      \flag{confidence: high on the temperature comparison; medium on the
      propagated effect on $f_{\rm ion}$, because I did not recompute
      $\alpha(T)$ or re-run the balance; and Dhandha et al.\ explicitly exclude
      flexible Pop~III source models, which is the regime Douna et al.\ occupy.}
      \textbf{This is the constraint your paper most needs to address head-on.}

\item \textbf{The second-strongest conflict is with the transport model.}
      Owen et al.\ (2018) on magnetic confinement, Jana \& Nath (2018) and
      Gessey-Jones et al.\ (2023) on the consequences of diffusion, and
      Yokoyama \& Ohira (2023) and Bera et al.\ (2023) on channels that only
      exist in a magnetized medium, jointly indicate that H6 --- deferring
      magnetic fields to a companion paper --- is the load-bearing simplification.
      \flag{confidence: medium-high} --- the papers agree that B matters and
      \emph{disagree with each other} on whether the result is patchy or
      uniform heating.

\item \textbf{The ionization claim is the most exposed result.}
      Leite et al.\ (2017) obtain a rate of the same order as Douna et al.\ and
      call it negligible. Nothing published since has revisited this.
      \flag{confidence: high that the disagreement in interpretation is real;
      low that it is resolvable from the two papers alone.}

\item \textbf{Carrier mismatch is pervasive and is the main weakness of this
      audit.} Ten of the thirteen entries model CR \emph{protons} or X-ray
      photons. Douna et al.\ model CR \emph{electrons}. Arguments that transfer
      cleanly are the geometric/transport ones (rigidity-based confinement,
      deposition range) and the observational ones (IGM temperature). Arguments
      about energy budget and spectrum do not transfer without a calculation
      neither paper performs. \flag{This is stated as a limitation of the
      literature, not of the search: no paper modelling CR-electron transport
      from high-$z$ microquasars other than Douna et al.\ (2018) and Tueros et
      al.\ (2014) was found.}
\end{enumerate}

\subsection{Explicitly not verified}

\begin{itemize}[leftmargin=*,itemsep=2pt]
\item No numerical quantity from any paper was independently recomputed. Every
      number quoted is as printed by its authors. No cross-check with
      \texttt{sympy} or by limiting case was performed, because no derivation
      was reproduced. \flag{no cross-check performed}
\item The comparison of HERA's $T_K$ bound at $z=10.4$ with Douna et al.'s
      assumption at $z=10$ ignores the (small, monotonic) redshift difference;
      the direction of the inequality is unaffected but the factor is
      approximate.
\item Search recall \emph{was} validated against NASA ADS (Section~\ref{sec:adsdelta}),
      which added four citing papers and two conflicting ones. ADS full-text
      search (\texttt{full:}) covers only the fraction of the corpus for which
      ADS holds full text; a conflicting statement buried in the body of a paper
      outside that fraction would still be missed.
\item Two papers (Fornasini et al.\ 2019, 2020) were assessed from their ADS
      abstracts only and judged supporting. The Sazonov \& Khabibullin (2017)
      reversal in Section~\ref{sec:screened} shows abstract-level screening can
      invert on full reading, so this is a real residual risk.
\item Excerpts were extracted mechanically and hand-checked against the printed
      form for mathematical symbols only. Long quotations were not re-read
      against the published (journal) versions; all are from the arXiv PDFs
      stored in \texttt{against/}.
\end{itemize}

% =====================================================================
\section*{Running reference list}
\addcontentsline{toc}{section}{Running reference list}
% =====================================================================

\begin{itemize}[leftmargin=*,itemsep=1pt,label={}]\small
\item Abdurashidova, Z., et al.\ (HERA Collaboration) 2023, ApJ 945, 124. DOI 10.3847/1538-4357/acaf50. arXiv:2210.04912.
\item Bera, A., Samui, S.\ \& Datta, K.~K.\ 2023, MNRAS 519, 4869. DOI 10.1093/mnras/stac3814. arXiv:2202.12308.
\item Bosch-Ramon, V.\ 2018, A\&A 617, L3. DOI 10.1051/0004-6361/201833952. arXiv:1808.08911.
\item Carvalho, L., Escobar, G.~J.\ \& Pellizza, L.~J.\ 2024, BAAA 65, 240. Bibcode 2024BAAA...65..240C.
\item Dasgupta, S., Liu, B., Fialkov, A., et al.\ 2026, arXiv:2604.11542.
\item Dhandha, J., Fialkov, A., Gessey-Jones, T., et al.\ 2025, MNRAS 544, 1608. DOI 10.1093/mnras/staf1736. arXiv:2508.13761.
\item Escobar, G.~J., Pellizza, L.~J.\ \& Romero, G.~E.\ 2021, BAAA 62, 265. Bibcode 2021BAAA...62..265E.
\item Fornasini, F.~M., Kriek, M., Sanders, R.~L., et al.\ 2019, ApJ 885, 65. DOI 10.3847/1538-4357/ab4653. arXiv:1909.08635.
\item Fornasini, F.~M., Civano, F.\ \& Suh, H.\ 2020, MNRAS 495, 771. DOI 10.1093/mnras/staa1211. arXiv:2004.13033.
\item Garate N\'u\~nez, L.~P., Escobar, G.~J., Pellizza, L.~J.\ \& Bosch-Ramon, V.\ 2021, BAAA 62, 234. Bibcode 2021BAAA...62..234G.
\item Douna, V.~M., Pellizza, L.~J., Laurent, P.\ \& Mirabel, I.~F.\ 2018, MNRAS 474, 3488. DOI 10.1093/mnras/stx2983. arXiv:1711.07374.
\item Escobar, G.~J., Pellizza, L.~J.\ \& Romero, G.~E.\ 2021, A\&A 650, A136. DOI 10.1051/0004-6361/202039860. arXiv:2104.11975.
\item Escobar, G.~J., Pellizza, L.~J.\ \& Romero, G.~E.\ 2022, A\&A 665, A61. DOI 10.1051/0004-6361/202142753. arXiv:2207.08633.
\item Gessey-Jones, T., Fialkov, A., de Lera Acedo, E., Handley, W.~J.\ \& Barkana, R.\ 2023, MNRAS 526, 4262. DOI 10.1093/mnras/stad3014. arXiv:2304.07201.
\item Jana, R.\ \& Nath, B.~B.\ 2018, MNRAS 479, 153. DOI 10.1093/mnras/sty1481. arXiv:1806.01295.
\item Kaur, H.~D., Qin, Y.\ \& Mesinger, A.\ 2022, MNRAS 513, 5097. DOI 10.1093/mnras/stac1226. arXiv:2203.10851.
\item Lehmer, B.~D., Eufrasio, R.~T., Basu-Zych, A., et al.\ 2021, ApJ 907, 17. DOI 10.3847/1538-4357/abcec1. arXiv:2011.09476.
\item Lehmer, B.~D., Monson, E.~B., Eufrasio, R.~T., et al.\ 2024, arXiv:2410.19901.
\item Leite, N., Evoli, C., D'Angelo, M., Ciardi, B., Sigl, G.\ \& Ferrara, A.\ 2017, MNRAS 469, 416. DOI 10.1093/mnras/stx805. arXiv:1703.09337.
\item Leong, K.-H.\ \& Meiksin, A.\ 2025, arXiv:2509.10255.
\item Mirabel, I.~F.\ \& Rodr\'iguez, L.~F.\ 2022, New Astron.\ Rev.\ 94, 101642. DOI 10.1016/j.newar.2022.101642. arXiv:2203.12741.
\item Owen, E.~R., Jacobsen, I.~B., Wu, K.\ \& Surajbali, P.\ 2018, MNRAS 481, 666. DOI 10.1093/mnras/sty2279. arXiv:1808.07837.
\item Owen, E.~R., Jin, X., Wu, K.\ \& Ferreras, I.\ 2019, MNRAS 484, 1645. DOI 10.1093/mnras/stz060. arXiv:1901.01411.
\item Pellizza, L.~J., Badaracco, M.~B., Carvalho, L., Escobar, G.~J.\ \& Bignone, L.~A.\ 2025, BAAA 66, 377. Bibcode 2025BAAA...66..377P.
\item Sartorio, N.~S., Fialkov, A., Hartwig, T., et al.\ 2023, MNRAS 521, 4039. DOI 10.1093/mnras/stad697. arXiv:2303.03435.
\item Sazonov, S.~Yu.\ \& Khabibullin, I.~I.\ 2017, Astron.\ Lett.\ 43, 211. DOI 10.1134/S1063773717040077. arXiv:1612.01262.
\item Sazonov, S.\ \& Khabibullin, I.\ 2018, MNRAS 476, 2530. DOI 10.1093/mnras/sty442. arXiv:1712.04831.
\item Singh, S., Nambissan T., J., Subrahmanyan, R., et al.\ 2022, Nature Astronomy 6, 607. DOI 10.1038/s41550-022-01610-5. arXiv:2112.06778.
\item Sotomayor Checa, P., Romero, G.~E.\ \& Bosch-Ramon, V.\ 2021, Ap\&SS 366, 1. DOI 10.1007/s10509-020-03911-5. arXiv:2012.11664.
\item Tueros, M., del Valle, M.~V.\ \& Romero, G.~E.\ 2014, A\&A 570, L3. arXiv:1409.6225.
\item Yang, Y., Long, H.\ \& Hirata, C.~M.\ 2024, ApJ 965, 111. DOI 10.3847/1538-4357/ad3237. arXiv:2311.18721.
\item Yokoyama, S.~L.\ \& Ohira, Y.\ 2023, MNRAS 523, 3671. DOI 10.1093/mnras/stad1596. arXiv:2302.07028.
\end{itemize}

\end{document}
